\section{省级政治、公共信息与帝国转换的证据}
1909年至1912年的转型同时发生在议政机构、公司与铁路、军队、报刊、电报、地方财政和谈判桌上。本节以地方政治与公共信息为中心重读账本，并明确指出：女性、家庭成员、基层雇工与流动者较少直接进入制度档案，不能让其缺席变成政治变化没有触及他们的证明。
\subsection{1905—1909：改革行动、地方组织与材料边界}
这些节点按账本日期相连，却不构成一条预定的现代化道路。每一项制度、项目或冲突都需要在其财政来源、人员、区域和信息传播条件中重新定位。
\hypertarget{public:EQ-1909-PROVINCIAL-ASSEMBLIES}{}\paragraph{宣统元年（1909年）}各省咨议局开办，地方精英获得新的公开政治场域。它应放入制度语言转为省级组织、军队调动、报刊消息与谈判实践的过程中。咨议、铁路、独立或退位各有不同的行动者，不能因最终结果而被预设为同一政治意图。 核心取证为《宣统政纪》宣统元年相关条；宫中朱批奏折、内阁大库题本与《清史稿》相关志传。上谕、题本与部议可确认规则如何被提出和机构如何分工；它们不能越过传递、落实与反馈的断点。可继续追问的研究问题是宫廷政治、制度运作与中央—地方信息链研究。账本列出的条件为预备立宪需要有限度吸纳地方意见。需要分层验证的影响为选举资格、代表性和实际权力均受限制。在性别与家庭史的尺度上，除非账簿、契约、学校名册、诉讼或私人书信给出可识别的当事人，不能从男性官员、军人或报人留下的叙述推演所有家庭成员的选择。来源保留：以宣统元年（1909年）的同期文书为定位起点；官修记录、奏报或外交文本服务特定机构立场，具体日期、规模、地方执行与对手视角须以独立材料复核。
\hypertarget{public:EQ-1909-RAILWAY-RECOVERY-MOVEMENT}{}\paragraph{宣统元年（1909年）}铁路国有化和路权问题引发各省商绅、公司和官府的争论。它应放入制度语言转为省级组织、军队调动、报刊消息与谈判实践的过程中。咨议、铁路、独立或退位各有不同的行动者，不能因最终结果而被预设为同一政治意图。 核心取证为《宣统政纪》宣统元年相关条；户部奏销、盐政、河工或海关档案。上谕、题本与部议可确认规则如何被提出和机构如何分工；它们不能越过传递、落实与反馈的断点。可继续追问的研究问题是财政账册、市场网络、灾害环境与区域经济研究。账本列出的条件为外债、地方集资与中央财政需求相冲突。需要分层验证的影响为“收回路权”诉求包含不同地区和利益群体。在性别与家庭史的尺度上，除非账簿、契约、学校名册、诉讼或私人书信给出可识别的当事人，不能从男性官员、军人或报人留下的叙述推演所有家庭成员的选择。来源保留：以宣统元年（1909年）的同期文书为定位起点；官修记录、奏报或外交文本服务特定机构立场，具体日期、规模、地方执行与对手视角须以独立材料复核。
\subsection{1910—1914：改革行动、地方组织与材料边界}
这些节点按账本日期相连，却不构成一条预定的现代化道路。每一项制度、项目或冲突都需要在其财政来源、人员、区域和信息传播条件中重新定位。
\hypertarget{public:EQ-1910-ZIZHENG-YUAN}{}\paragraph{宣统二年（1910年）}资政院开院，中央层面的立宪咨询机构开始运作。它应放入制度语言转为省级组织、军队调动、报刊消息与谈判实践的过程中。咨议、铁路、独立或退位各有不同的行动者，不能因最终结果而被预设为同一政治意图。 核心取证为《宣统政纪》宣统二年相关条；宫中朱批奏折、内阁大库题本与《清史稿》相关志传。上谕、题本与部议可确认规则如何被提出和机构如何分工；它们不能越过传递、落实与反馈的断点。可继续追问的研究问题是宫廷政治、制度运作与中央—地方信息链研究。账本列出的条件为朝廷试图展示宪政进程并控制代表参与。需要分层验证的影响为咨询权与立法权的界限成为持续争议。在性别与家庭史的尺度上，除非账簿、契约、学校名册、诉讼或私人书信给出可识别的当事人，不能从男性官员、军人或报人留下的叙述推演所有家庭成员的选择。来源保留：以宣统二年（1910年）的同期文书为定位起点；官修记录、奏报或外交文本服务特定机构立场，具体日期、规模、地方执行与对手视角须以独立材料复核。
\hypertarget{public:EQ-1910-FAMINE-AND-UNREST}{}\paragraph{宣统二年（1910年）}灾荒、物价和财政压力持续影响晚清地方社会与政治动员。它应放入制度语言转为省级组织、军队调动、报刊消息与谈判实践的过程中。咨议、铁路、独立或退位各有不同的行动者，不能因最终结果而被预设为同一政治意图。 核心取证为《宣统政纪》宣统二年相关条；地方志、督抚奏折及地方档案。实录的时间骨架可定位国家所承认的节点，却须同省档、报刊、私人日记与地方记录核验实际的消息速度和覆盖范围。可继续追问的研究问题是地方档案、迁徙、宗教、族群与社会动员研究。账本列出的条件为自然环境、市场波动和行政能力交互作用。需要分层验证的影响为地方灾情与政治行动不可简单建立单因果关系。在性别与家庭史的尺度上，除非账簿、契约、学校名册、诉讼或私人书信给出可识别的当事人，不能从男性官员、军人或报人留下的叙述推演所有家庭成员的选择。来源保留：以宣统二年（1910年）的同期文书为定位起点；官修记录、奏报或外交文本服务特定机构立场，具体日期、规模、地方执行与对手视角须以独立材料复核。
\hypertarget{public:EQ-1911-RAILWAY-NATIONALIZATION}{}\paragraph{宣统三年（1911年）}清廷宣布铁路干线国有化并以外债筹资，保路运动迅速扩大。它应放入制度语言转为省级组织、军队调动、报刊消息与谈判实践的过程中。咨议、铁路、独立或退位各有不同的行动者，不能因最终结果而被预设为同一政治意图。 核心取证为《宣统政纪》宣统三年相关条；宫中朱批奏折、内阁大库题本与《清史稿》相关志传。上谕、题本与部议可确认规则如何被提出和机构如何分工；它们不能越过传递、落实与反馈的断点。可继续追问的研究问题是宫廷政治、制度运作与中央—地方信息链研究。账本列出的条件为中央财政、外债与地方商办铁路利益冲突。需要分层验证的影响为政策公告、地方抗议和后续镇压应分开记录。在性别与家庭史的尺度上，除非账簿、契约、学校名册、诉讼或私人书信给出可识别的当事人，不能从男性官员、军人或报人留下的叙述推演所有家庭成员的选择。来源保留：以宣统三年（1911年）的同期文书为定位起点；官修记录、奏报或外交文本服务特定机构立场，具体日期、规模、地方执行与对手视角须以独立材料复核。
\hypertarget{public:EQ-1911-SICHUAN-RAILWAY-PROTEST}{}\paragraph{宣统三年（1911年）}四川保路运动升级，赵尔丰处置引发更大政治危机。它应放入制度语言转为省级组织、军队调动、报刊消息与谈判实践的过程中。咨议、铁路、独立或退位各有不同的行动者，不能因最终结果而被预设为同一政治意图。 核心取证为《宣统政纪》宣统三年相关条；地方志、督抚奏折及地方档案。实录的时间骨架可定位国家所承认的节点，却须同省档、报刊、私人日记与地方记录核验实际的消息速度和覆盖范围。可继续追问的研究问题是地方档案、迁徙、宗教、族群与社会动员研究。账本列出的条件为地方股东、商绅、学生和官府对路权与财政有不同诉求。需要分层验证的影响为参与规模和暴力过程须依当地档案与报刊核验。在性别与家庭史的尺度上，除非账簿、契约、学校名册、诉讼或私人书信给出可识别的当事人，不能从男性官员、军人或报人留下的叙述推演所有家庭成员的选择。来源保留：以宣统三年（1911年）的同期文书为定位起点；官修记录、奏报或外交文本服务特定机构立场，具体日期、规模、地方执行与对手视角须以独立材料复核。
\hypertarget{public:EQ-1911-WUCHANG-UPRISING}{}\paragraph{宣统三年（1911年）}武昌起义爆发，新军与革命团体控制武汉部分地区。它应放入制度语言转为省级组织、军队调动、报刊消息与谈判实践的过程中。咨议、铁路、独立或退位各有不同的行动者，不能因最终结果而被预设为同一政治意图。 核心取证为《宣统政纪》宣统三年相关条；军机处档、战事方略及地方奏报。实录的时间骨架可定位国家所承认的节点，却须同省档、报刊、私人日记与地方记录核验实际的消息速度和覆盖范围。可继续追问的研究问题是战争动员、军需、战场暴力与地方社会研究。账本列出的条件为新军组织、革命网络和铁路危机相互叠加。需要分层验证的影响为起义的即时参与者、城市控制与后续省份响应需分层叙述。在性别与家庭史的尺度上，除非账簿、契约、学校名册、诉讼或私人书信给出可识别的当事人，不能从男性官员、军人或报人留下的叙述推演所有家庭成员的选择。来源保留：以宣统三年（1911年）的同期文书为定位起点；官修记录、奏报或外交文本服务特定机构立场，具体日期、规模、地方执行与对手视角须以独立材料复核。
\hypertarget{public:EQ-1911-PROVINCIAL-INDEPENDENCE}{}\paragraph{宣统三年（1911年）}多省宣布脱离清廷，帝国行政和财政网络快速分裂。它应放入制度语言转为省级组织、军队调动、报刊消息与谈判实践的过程中。咨议、铁路、独立或退位各有不同的行动者，不能因最终结果而被预设为同一政治意图。 核心取证为《宣统政纪》宣统三年相关条；宫中朱批奏折、内阁大库题本与《清史稿》相关志传。上谕、题本与部议可确认规则如何被提出和机构如何分工；它们不能越过传递、落实与反馈的断点。可继续追问的研究问题是宫廷政治、制度运作与中央—地方信息链研究。账本列出的条件为武昌起义、地方军政力量和立宪派选择共同推动变化。需要分层验证的影响为“独立”声明不等于各省形成稳定统一控制。在性别与家庭史的尺度上，除非账簿、契约、学校名册、诉讼或私人书信给出可识别的当事人，不能从男性官员、军人或报人留下的叙述推演所有家庭成员的选择。来源保留：以宣统三年（1911年）的同期文书为定位起点；官修记录、奏报或外交文本服务特定机构立场，具体日期、规模、地方执行与对手视角须以独立材料复核。
\hypertarget{public:EQ-1911-YUAN-SHIKAI-RETURN}{}\paragraph{宣统三年（1911年）}清廷起用袁世凯处理北方军政和议和事务。它应放入制度语言转为省级组织、军队调动、报刊消息与谈判实践的过程中。咨议、铁路、独立或退位各有不同的行动者，不能因最终结果而被预设为同一政治意图。 核心取证为《宣统政纪》宣统三年相关条；宫中朱批奏折、内阁大库题本与《清史稿》相关志传。上谕、题本与部议可确认规则如何被提出和机构如何分工；它们不能越过传递、落实与反馈的断点。可继续追问的研究问题是宫廷政治、制度运作与中央—地方信息链研究。账本列出的条件为中央缺乏可依赖的军队与谈判人选。需要分层验证的影响为袁在清廷、北洋与革命派之间的策略塑造退位进程。在性别与家庭史的尺度上，除非账簿、契约、学校名册、诉讼或私人书信给出可识别的当事人，不能从男性官员、军人或报人留下的叙述推演所有家庭成员的选择。来源保留：以宣统三年（1911年）的同期文书为定位起点；官修记录、奏报或外交文本服务特定机构立场，具体日期、规模、地方执行与对手视角须以独立材料复核。
\hypertarget{public:EQ-1912-REPUBLIC-FOUNDING}{}\paragraph{宣统四年（1912年）}中华民国临时政府在南京成立，与清廷退位谈判并行。它应放入制度语言转为省级组织、军队调动、报刊消息与谈判实践的过程中。咨议、铁路、独立或退位各有不同的行动者，不能因最终结果而被预设为同一政治意图。 核心取证为《宣统政纪》宣统四年相关条；宫中朱批奏折、内阁大库题本与《清史稿》相关志传。上谕、题本与部议可确认规则如何被提出和机构如何分工；它们不能越过传递、落实与反馈的断点。可继续追问的研究问题是宫廷政治、制度运作与中央—地方信息链研究。账本列出的条件为各省独立与南北议和创造新的政权竞争格局。需要分层验证的影响为临时政府的法理主张与实际控制范围需要区分。在性别与家庭史的尺度上，除非账簿、契约、学校名册、诉讼或私人书信给出可识别的当事人，不能从男性官员、军人或报人留下的叙述推演所有家庭成员的选择。来源保留：以宣统四年（1912年）的同期文书为定位起点；官修记录、奏报或外交文本服务特定机构立场，具体日期、规模、地方执行与对手视角须以独立材料复核。
\hypertarget{public:EQ-1912-QING-ABDICATION}{}\paragraph{宣统四年（1912年）}宣统帝颁布退位诏书，清帝国终结。它应放入制度语言转为省级组织、军队调动、报刊消息与谈判实践的过程中。咨议、铁路、独立或退位各有不同的行动者，不能因最终结果而被预设为同一政治意图。 核心取证为《宣统政纪》宣统四年相关条；宫中朱批奏折、内阁大库题本与《清史稿》相关志传。上谕、题本与部议可确认规则如何被提出和机构如何分工；它们不能越过传递、落实与反馈的断点。可继续追问的研究问题是宫廷政治、制度运作与中央—地方信息链研究。账本列出的条件为南北谈判、财政军政压力与皇室待遇安排共同促成退位。需要分层验证的影响为退位文本规定的优待与地方政治现实须分别追踪。在性别与家庭史的尺度上，除非账簿、契约、学校名册、诉讼或私人书信给出可识别的当事人，不能从男性官员、军人或报人留下的叙述推演所有家庭成员的选择。来源保留：以宣统四年（1912年）的同期文书为定位起点；官修记录、奏报或外交文本服务特定机构立场，具体日期、规模、地方执行与对手视角须以独立材料复核。
